% 无向边与一对反向弧：两者都允许在 u,v 之间往返，但后者是两个独立对象。
\begin{tikzpicture}[panelmark/.style={font=\figurefont\fontsize{8}{10}\selectfont,text=FigureInk,anchor=north,inner sep=0pt,outer sep=0pt},
>=Stealth,
  v/.style={circle,draw=FigureStroke,fill=FigureBlueFill,line width=.8pt,minimum size=6.5mm,inner sep=0pt},
  e/.style={draw=FigureStroke,line width=.8pt},
  a/.style={draw=FigureStroke,line width=.8pt,->,shorten >=1pt}]
\begin{scope}[local bounding box=panelrow]
\begin{scope}[local bounding box=panela]
  \node[v] (u) at (0,0) {$u$}; \node[v] (w) at (1.6,0) {$v$};
  \draw[e] (u)--(w);

\end{scope}
\begin{scope}[xshift=3.4cm,local bounding box=panelb]
  \node[v] (u) at (0,0) {$u$}; \node[v] (w) at (1.6,0) {$v$};
  \draw[a] (u) to[bend left=25] (w);
  \draw[a] (w) to[bend left=25] (u);

\end{scope}
\end{scope}
% Subfigure marks: 8 pt, centred below each panel; common top 3 mm below this row.
\coordinate (panelrowbase) at ($(panelrow.south)+(0,-3mm)$);
\node[panelmark] at (panela.center |- panelrowbase) {(a)};
\node[panelmark] at (panelb.center |- panelrowbase) {(b)};
\end{tikzpicture}
